\chapter{Teorias de primeira ordem}
A Lógica de Primeira Ordem é um sistema formal que permite expressar propriedades e relações entre objetos por meio dos quantificadores ``para todo'' ($\forall$) e ``existe'' ($\exists$), além de símbolos para relações, funções e constantes.

A Teoria dos Conjuntos é uma \emph{teoria de primeira ordem}\index{teoria de primeira ordem}, isto é, tem como base a \emph{lógica de primeira ordem}\index{lógica de primeira ordem}. 
Assim, antes de adentrarmos a teoria dos conjuntos, discutiremos noções fundamentais sobre linguagens de primeira ordem e sobre a própria lógica.

Este capítulo não pretende, de forma alguma, ser uma introdução completa à lógica de primeira ordem. 
Nosso objetivo é apenas fixar a linguagem e tornar o texto autocontido, fornecendo ao leitor as ferramentas necessárias para a compreensão dos capítulos seguintes.
Em apêndices futuros, não essenciais à progressão do texto, apresentaremos aprofundamentos adicionais.

\subimport{}{chap1-polish}
\subimport{}{chap1-firstorder}
\subimport{}{chap1-syntax}
\subimport{}{chap1-tautologies}
\subimport{}{chap1-theories}